% Artifact Appendix template for the NDSS Artifact Evaluation
% version 1.1 (20250525)

% Remove the following block when merging the appendix with the camera-ready paper.
%%%
\documentclass[conference]{IEEEtran}
\usepackage[most]{tcolorbox}
\usepackage{listings}
\usepackage[scaled=0.88]{newtxtt}
\usepackage[most]{tcolorbox}
\usepackage{listings}

\newtcblisting{commandbox}{
  listing only,
  listing options={
    basicstyle=\ttfamily,
    breaklines=true,
    columns=fullflexible,
  },
  colback=black!2,
  colframe=black!30,
  boxrule=0pt,
  leftrule=1.2pt,
  arc=0pt,
  left=6pt,
  right=4pt,
  top=3pt,
  bottom=3pt,
  boxsep=0pt,
  before skip=5pt,
  after skip=5pt,
}
\pagestyle{plain}
\usepackage{url}
\usepackage{hyperref}
\begin{document}
%%%

\appendices
\section{Artifact Appendix}

This artifact accompanies \emph{Isotaxis: Optimal Asynchronous Byzantine
Consensus with Ordering Linearizability}.  It provides the Isotaxis and Pompe
implementations, a common HotStuff backend, four-replica Docker experiments,
the global EC2 workflow, and the data and script for Figs.~2, 3, and 5--7.
These components evaluate the latency, throughput--latency, and ordering-layer
denial-of-service (DoS) results in Section VIII of the paper.  Operational
details are in \texttt{README-AE.md}.

\subsection{Description \& Requirements}

The main Rust crate contains Isotaxis, Pompe, clients, and adversarial nodes;
the bundled second crate is their modified HotStuff backend.  The two Compose
files run the normal and attack cases.  Global deployment scripts are under
\path{ec2}, and \path{gen-fig.py} regenerates the plots.

\subsubsection{How to access}

The repository is available at
\url{https://github.com/xrxrxrxrxr/Isotaxis-ArtifactEvaluation}.  Clone the
evaluated revision:

\begin{commandbox}
git clone \
  https://github.com/xrxrxrxrxr/\
Isotaxis-ArtifactEvaluation
cd Isotaxis-ArtifactEvaluation
\end{commandbox}

The evaluated commit is recorded in \texttt{README-AE.md}; the camera-ready
appendix will cite the AEC-approved archive DOI.  The authors' code is
Apache-2.0; third-party components retain their original licenses.

\subsubsection{Hardware dependencies}

The local four-replica run needs an x86-64 host with 8 logical cores, 16~GiB
RAM, and 20~GiB free disk; no GPU is needed.  CPU-affinity settings refer to
cores 0--3.  Docker Desktop can run the functional test, but its measurements
are not directly comparable to the paper.  A full rerun needs one dedicated
EC2 \texttt{c6a.2xlarge} (8 vCPUs, 16~GiB) per process and a
\texttt{t3.micro} client, distributed over 11 regions for up to 100 processes.

\subsubsection{Software dependencies}

Local experiments require Git, Docker Engine, and Docker Compose; testing used
Engine 28.3.2, Compose v2.38.2, and Buildx v0.25.0.  Rust is built inside the
image.  Plotting needs Python~3.9+ and \texttt{matplotlib==3.9.4}, pinned in
\path{requirements-plot.txt}.  E3 additionally needs AWS CLI v2,
\texttt{jq}, SSH/SCP, GNU Make, credentials, and the EC2 fleet above.

\subsubsection{Benchmarks}

No dataset is required: the client generates synthetic transactions.
\texttt{ORDERING\_BATCH\_SIZE} is the nominal transactions per ordering input;
the paper uses one for sparse latency and sweeps up to $1.6\times10^{5}$ for
throughput.  Processed series for Figs.~2, 3, and 5--7 are embedded in
\path{gen-fig.py}; live runs produce client and per-replica logs.

\subsection{Artifact Installation \& Configuration}

From the repository root, verify the revision and both Compose configurations:

\begin{commandbox}
git rev-parse HEAD
docker compose config --quiet
docker compose --env-file .attack.env \
  -f docker-compose.attack.yml config -q
\end{commandbox}

The revision must match \texttt{README-AE.md}, and both configuration checks
must succeed.  E1 builds Rust inside Docker.  For E2:

\begin{commandbox}
python3 -m venv .venv-ae
source .venv-ae/bin/activate
python -m pip install -r \
  requirements-plot.txt
\end{commandbox}

\subsection{Experiment Workflow}

E0 validates the package; E1.1 and E1.2 run the local normal and DoS paths; E2
renders the archived paper data; and optional E3 collects new global EC2 data.
Because \texttt{run\_test.sh} clears \texttt{logs}, preserve each run first.
For comparisons, use at least three repetitions with the same commit,
environment, hardware, batch size, and attack rate, then compare medians.

\subsection{Major Claims}

\begin{itemize}
    \setlength{\itemsep}{0pt}
    \setlength{\parskip}{0pt}
    \setlength{\parsep}{0pt}
    \item \textbf{(C1) Sparse-load scalability.}  The confirmation latency of
    Pompe grows faster with replica count.  At $n=100$, Fig.~5 reports about
    22~s for Pompe and 12~s for Isotaxis.  E1.1 checks the path; E2 checks the
    archived result.

    \item \textbf{(C2) Throughput--latency trade-off.}  Batching initially
    raises throughput until resource limits; Isotaxis's advantage grows with
    $n$.  At $n=100$ and batch size $\approx1.3\times10^{4}$, it exceeds
    28,000 transactions/s, nearly twice Pompe.  Figs.~3 and~6 are checked by E2
    and optionally rerun by E3.

    \item \textbf{(C3) Robustness to ordering-layer DoS.}  With one of 30
    processes flooding, Fig.~7 shows Isotaxis changing from 2.597 to 3.067~s
    and from 1,155 to 978 transactions/s; Pompe changes from 2.030 to 33.474~s
    and from 1,478 to 90 transactions/s.  E1.2 checks the path and E2 the data.
\end{itemize}

The paper's safety, liveness, resilience, and asymptotic communication claims
are established analytically and are not treated as artifact-evaluation
claims.

\subsection{Evaluation}

\subsubsection{Experiment (E0): package and build validation}

\textit{[5--20 human-minutes plus first-build time].}
\textit{Preparation/Execution:} run the installation checks above.
\textit{Results:} all commands exit zero and the revision matches; E1 must
subsequently build four replica containers and the load generator.

\subsubsection{Experiment (E1.1): local normal execution}

\textit{[About 5 human-minutes and 4 compute-minutes per run].}
\textit{Preparation:} in \texttt{.env}, set
\texttt{CLIENT\_ORDERING\_MODE=smrol} and
\texttt{ORDERING\_BATCH\_SIZE=1} for the four-node sparse-latency path.
Use a larger batch for a throughput check.  \textit{Execution:}

\begin{commandbox}
./run_test.sh --rebuild
./run_test.sh --reuse
\end{commandbox}

Repeat with \texttt{CLIENT\_ORDERING\_MODE=pompe} for the baseline.  Preserve
the logs before switching.  \textit{Results:} \path{logs/load_test.log} contains
aggregate ordering/consensus latency, and
\path{logs/node/consensus-node*.log} shows continuing commits and throughput.
This is the four-node counterpart of Figs.~3, 5, and~6, not their global scale.

\subsubsection{Experiment (E1.2): local DoS execution}

\textit{[About 5 human-minutes and 4 compute-minutes per run].}
\textit{Preparation:} preserve matching E1.1 logs, then set in
\texttt{.attack.env}:

\begin{commandbox}
CLIENT_ORDERING_MODE=smrol
ADVERSARY_MODE=smrol
CLIENT_EXCLUDED_NODE_IDS=1
ATTACK_SLEEP_MICROS=100
\end{commandbox}

For Pompe, change both protocol variables to \texttt{pompe} and set
\texttt{POMPE\_ATTACK\_WORKERS=4}.  \textit{Execution:}

\begin{commandbox}
./run_test.sh adversary --rebuild
./run_test.sh adversary --reuse
\end{commandbox}

\textit{Results:} node1 reports the pure attacker and HotStuff participation;
all replicas keep committing and the client prints its aggregate report.  Pompe
also prints \texttt{[pompe-attacker][rate]}.  Across matched repetitions,
Pompe's median ordering latency should degrade markedly, while Isotaxis remains
comparatively stable, matching Fig.~7 qualitatively.  Illustrative four-node
logs are in \path{log-examples/attack-logs}; new results are in \path{logs}.


\subsubsection{Experiment (E2): regenerate the paper figures}

\textit{[About 2 human-minutes; less than 1 compute-minute].}
\textit{Preparation/Execution:}

\begin{commandbox}
source .venv-ae/bin/activate
python gen-fig.py
\end{commandbox}

\textit{Results:} ten PDFs appear under path \path{python-fig}: those figures are the
\path{attacks-intro}, \path{tradeoff-100-intro}, and \path{latency} PDFs map to
Figs.~2, 3, and~5; the six \path{tps-*}/\path{tradeoff-*} PDFs map to Fig.~6;
and \path{attacks-exp.pdf} maps to Fig.~7.  They must show the C1--C3 values
above.  E2 renders archived processed series; it is not a new AWS measurement.

\subsubsection{Experiment (E3): global AWS execution for Fig.~6}

\textit{[Optional; about 5 compute-minutes per point or 4.5 compute-hours for
the full sweep].}  This costly, state-changing workflow selects existing
instances; first inspect the plan and dry-run:

\begin{commandbox}
cd ec2
NODE_INSTANCE_TYPE=c6a.2xlarge \
  ./start_instances.sh \
  --plan plans/plan-16.json --dry-run
\end{commandbox}

Plans 16, 64, and 100 correspond to the Fig.~6 panels.  On an authorized
account with initialized instances, \textit{execute:}

\begin{commandbox}
./start_instances.sh \
  --plan plans/plan-16.json
NODE_IP_SOURCE=public ./get_ips.sh
\end{commandbox}

\path{get_ips.sh} invokes \texttt{make deploy}, which runs, collects logs, and
stops the instances; do not invoke \texttt{make deploy/start} again.  Repeat
the documented batch/load sweep for both protocols and record the plan and
configuration.  \textit{Results:} \path{ec2/ec2-logs} contains client latency
and committed throughput for reconstructing the selected Fig.~6 panel.

\subsection{Customization}

Use \texttt{.env} for normal runs and \texttt{.attack.env} for attacks.  Main
controls are the protocol, ordering batch size, excluded client targets, attack
interval, and Pompe worker count.  An \path{ec2/plans/plan-x.json} selects the
distributed replica count and regions.  Record any CPU-affinity changes.

% Remove the following block when merging the appendix with the camera-ready paper.
%%%
\end{document}
%%%
